% Fichier généré automatiquement par tools/generate_corrections.py.
% Source : DYN/DYN-04-TorseurDynamique/STOCK/07_RR3D_02/07_RR3D_02.tex
% Statut : complété (3/3 question(s)).
\begin{corrigebox}[Corrigé complété]
\CorrigeQuestion{1}
On choisit les inconnues et les conventions positives de la figure,
        on écrit les relations de conservation/fermeture adaptées, puis on résout le
        système obtenu. Le résultat symbolique doit être homogène et vérifié dans les cas
        limites (entrée nulle, rapport unitaire ou position de référence).
\CorrigeQuestion{2}
\CorrigeSource{Réponse reprise d'un exercice homonyme de la banque ; la provenance exacte est indiquée dans le commentaire du fichier.}
% Réponse reprise de : CIN/CIN-02-VitesseAcceleration/07_RR3D_02/07_RR3D_02.tex
$\torseurcin{V}{2}{0} = \torseurl{\dot{\theta}\vk{1}+ \dot{\varphi}\vi{1}}{\left(R+\ell\right)\dot{\theta}\vj{1} -r\dot{\theta}\cos\varphi \vi{1} + r\dot{\varphi}\vk{2}}{C}$
\CorrigeQuestion{3}
\CorrigeSource{Réponse reprise d'un exercice homonyme de la banque ; la provenance exacte est indiquée dans le commentaire du fichier.}
% Réponse reprise de : CIN/CIN-02-VitesseAcceleration/07_RR3D_02/07_RR3D_02.tex
$\vectg{C}{2}{0} = \deriv{\vectv{C}{2}{0}}{\rep{0}}$
 
 $ = \deriv{\left(R+\ell\right)\dot{\theta}\vj{1} -r\dot{\theta}\cos\varphi \vi{1} + r\dot{\varphi}\vk{2}}{\rep{0}}$
 
 Calculons : 
\begin{itemize}
\item $\deriv{\vi{1}}{\rep{0}}$ $=\vecto{1}{0} \wedge \vi{1}$ $=\dot{\theta}\vk{0} \wedge \vi{1}$ $=\dot{\theta}\vj{1}$.
\item $\deriv{\vj{1}}{\rep{0}}$ $=\vecto{1}{0} \wedge \vj{1}$ $=\dot{\theta}\vk{0} \wedge \vj{1}$ $=-\dot{\theta}\vi{1}$.
\item $\deriv{\vk{2}}{\rep{0}}$ $=\vecto{2}{0} \wedge \vk{2}$ 
$=\left(\dot{\theta}\vk{0}+ \dot{\varphi}\vi{1} \right) \wedge \vk{2}$
$=\dot{\theta}\vk{1} \wedge \vk{2}+ \dot{\varphi}\vi{2}  \wedge \vk{2}$
$=\dot{\theta}\sin\varphi\vi{1}- \dot{\varphi}\vj{2}$.
\end{itemize}

$\vectg{C}{2}{0}$ 
$ =\left(R+\ell\right)\ddot{\theta}\vj{1} -\left(R+\ell\right)\dot{\theta}^2\vi{1} 
-r\ddot{\theta}\cos\varphi \vi{1} +r\dot{\theta}\dot{\varphi}\sin\varphi \vi{1} -r\dot{\theta}^2\cos\varphi \vj{1} 
+ r\ddot{\varphi}\vk{2}+ r\dot{\varphi}\left(\dot{\theta}\sin\varphi\vi{1}- \dot{\varphi}\vj{2} \right)$.
\end{corrigebox}
